\documentclass[../../../main]{subfiles}

\ifSubfilesClassLoaded{%
    \zexternaldocument*{../../1.Category-Theory/exported/aux/main}%
    \zexternaldocument*{../../2.Homological-Algebra/exported/aux/main}%
    \zexternaldocument*{../../3.Sheaves/1.Sheaves/exported/aux/main}%
    \zexternaldocument*{1.Sheaves-of-modules}%
}{}

\begin{document}
\section{Ringed spaces} \label{sec:ringed-spaces}

\begin{definition}[Ringed space] \label{def:ringed-space}
  A \textbf{ringed space} is a pair \((X, \sheaf{O}_{X})\) consisting of a topological space \(X\) and a \hyperref[def:sheaf]{sheaf} of rings \(\sheaf{O}_{X}\) over \(X\).

  A \textbf{morphism} of ringed spaces \((X, \sheaf{O}_{X}) \to (Y, \sheaf{O}_{Y})\) is a pair \((f, f^{\#})\) where \(f: X \to Y\) is a continuous map and \(f^{\#}: \sheaf{O}_{y} \to f_{\ast}\sheaf{O}_{X}\) is a morphism of \hyperref[def:sheaf]{sheaves}.
\end{definition}

\subsection{Ringed spaces and continuous functions}

We wish to translate what we did for \hyperref[def:sheaf]{sheaves} in \autoref{ssec:presheaves-and-continuous-maps} to \hyperref[def:sheaf-of-modules]{sheaves of \(\sheaf{O}_{X}\)-modules}.

The direct image is direct.
\begin{definition} \label{def:direct-image-ox-mod}
  Let \((f, f^{\#}): (X, \sheaf{O}_{X}) \to (Y, \sheaf{O}_{Y})\) be a \hyperref[def:ringed-space]{morphism of ringed spaces}.
  Let \(\sheaf{F}\) be a \hyperref[def:sheaf-of-modules]{sheaf of \(\sheaf{O}_{X}\)-modules}.

  We define the \textbf{direct image} of \(\sheaf{F}\) via \(f\) as the \hyperref[def:sheaf-of-modules]{sheaf of \(\sheaf{O}_{Y}\)-modules} which has the \hyperref[def:direct-image-sheaf]{direct image sheaf} \(f_{\ast}\sheaf{F}\) as \hyperref[def:sheaf]{sheaf} of abelian groups and with module structure given by the map \(f^{\#}: \sheaf{O}_{Y} \to f_{\ast}\sheaf{O}_{X}\).

  This defines a \hyperref[def:functor]{functor}
  \begin{equation*}
    f_{\ast}: \Mod(\sheaf{O}_{X}) \to \Mod(\sheaf{O}_{Y})
  \end{equation*}
\end{definition}

However, as in the case for the \hyperref[def:inverse-image-sheaf]{inverse image for sheaves}, the inverse image of a \hyperref[def:sheaf-of-modules]{sheaf of \(\sheaf{O}_{X}\)-modules} needs more work.

\todo[inline]{Explain why \(\otimes\)}

\begin{definition} \label{def:inverse-image-ox-mod}
  Let \((f, f^{\#}): (X, \sheaf{O}_{X}) \to (Y, \sheaf{O}_{Y})\) be a \hyperref[def:ringed-space]{morphism of ringed spaces}.
  Let \(\sheaf{G}\) be an \(\Gamma(X, \sheaf{O}_{X})\)-module, where

  We define the \textbf{inverse image} of \(\sheaf{G}\) via \(f\) to be the \hyperref[def:sheaf-of-modules]{sheaf of \(\sheaf{O}_{X}\)-modules}
  \begin{equation*}
    f^{\ast}\sheaf{G} \coloneq \sheaf{O}_{X} \otimes_{f^{-1}\sheaf{O}_{Y}} f^{-1}\sheaf{G}
  \end{equation*}
  where \(f^{-1}\) is the \hyperref[def:inverse-image-sheaf]{inverse image sheaf}.

  This defines a \hyperref[def:functor]{functor}
  \begin{equation*}
    f^{\ast}: \Mod(\sheaf{O}_{Y}) \to \Mod(\sheaf{O}_{X})
  \end{equation*}

\end{definition}


\begin{theorem} \label{thm:direct-inverse-adjunction-ringed-space}
  Let \((f, f^{\#}): (X, \sheaf{O}_{X}) \to (Y, \sheaf{O}_{Y})\) be a \hyperref[def:ringed-space]{morphism of ringed spaces}.
  The \hyperref[def:functor]{functors} \(f_{\ast}\) and \(f^{\ast}\) are \hyperref[def:adjoint-functors]{adjoint}.

  In formula, if \(\sheaf{F}\) is a \hyperref[def:sheaf-of-modules]{sheaf of \(\sheaf{O}_{X}\)-modules} and \(\sheaf{G}\) is a \hyperref[def:sheaf-of-modules]{sheaf of \(\sheaf{O}_{Y}\)-modules} we have
  \begin{equation*}
    \Hom_{\Mod(\sheaf{O}_{X})} \left( f^{\ast}\sheaf{G}, \sheaf{F} \right) \cong \Hom_{\Mod(\sheaf{O}_{Y})} \left( \sheaf{G}, f_{\ast}\sheaf{F} \right)
  \end{equation*}
\end{theorem}
\begin{proof}
\begin{align*}
  \Hom_{\Mod (\sheaf{O}_{X})} \left( f^{\ast}\sheaf{G}, \sheaf{F} \right) &\cong \Hom_{\Mod (\sheaf{O}_{X})} \left( \sheaf{O}_{X} \otimes_{f^{-1} \sheaf{O}_{Y}} f^{-1}\sheaf{G}, \sheaf{F} \right) \\
                                                                       &\cong \Hom_{\Mod (f^{-1} \sheaf{O}_{Y})} \left( f^{-1}\sheaf{G}, \sheaf{F}_{f^{-1}\sheaf{O}_{Y}} \right) \\
                                                                       &\cong  \Hom_{\Mod(\sheaf{O}_{Y})} \left( \sheaf{G}, f_{\ast}\sheaf{F} \right)
\end{align*}
\end{proof}

With the \hyperref[def:adjoint-functors]{adjunction} between \(f_{\ast}\) and \(f^{\ast}\) out of the way, let's now ...

Let \((X, \sheaf{O}_{X})\) be a \hyperref[def:ringed-space]{ringed space} and \(p \in X\).

Consider the one-point topological space \(\{\ast\}\).
Recall that by \autoref{remark:empty-set-sheaf} we have that \(\Sh_{\{\ast\}}(\Ring) \cong \Ring\).
Therefore to make \(\{\ast\}\) into a \hyperref[def:ringed-space]{ringed space} we only need to \textit{accompany} it by a ring.

We are interested in the \hyperref[def:direct-image-ox-mod]{direct} and \hyperref[def:inverse-image-ox-mod]{inverse} images of the \hyperref[def:ringed-space]{morphisms of ringed spaces} induced by the maps
\begin{equation*}
  \pi: X \to \{\ast\} \qquad \iota_{p}: \{\ast\} \to X
\end{equation*}


Note that
\begin{align*}
  \pi^{\#}:& \sheaf{O}_{\pi} \to \pi_{\ast}\sheaf{O}_{X} = \Gamma(X, \sheaf{O}_{X}) \\
  \iota_{p}^{\#}:& \sheaf{O}_{X} \to \iota_{p, \ast}\sheaf{O}_{\iota_{p}} = \sheaf{O}_{\iota_{p}, p}
\end{align*}
Therefore we choose \(\sheaf{O}_{\pi} = \Gamma(X, \sheaf{O}_{X})\) and \(\sheaf{O}_{\iota_{p}} = \sheaf{O}_{X, p}\).

\begin{definition}[Global sections] \label{def:global-sections-ox-mod}
  Let \((X, \sheaf{O}_{X})\) be a \hyperref[def:ringed-space]{ringed space}.
  Let \(\sheaf{F}\) be a \hyperref[def:sheaf-of-modules]{sheaf of \(\sheaf{O}_{X}\)-modules}.
  Consider the \hyperref[def:ringed-space]{morphism of ringed spaces}
  \begin{equation*}
    (\pi, \pi^{\ast}): (X, \sheaf{O}_{X}) \to (\{\ast\}, \Gamma(X, \sheaf{O}_{X}))
  \end{equation*}

  
  We define the \textbf{global sections} of \(\sheaf{F}\) as the \hyperref[def:direct-image-ox-mod]{direct image} of \(\sheaf{F}\) via \(\pi\):
  \begin{equation*}
    \Gamma(X, \sheaf{F}) \coloneq \pi_{\ast}\sheaf{F}
  \end{equation*}

  Note that this defines a \hyperref[def:functor]{functor}
  \begin{equation*}
    \Gamma(X, -): \Mod(\sheaf{O}_{X}) \to \Mod(\Gamma(X, \sheaf{O}_{X}))
  \end{equation*}
\end{definition}

\begin{definition}["Constant" sheaf] \label{def:constant-sheaf-ox-mod}
  Let \((X, \sheaf{O}_{X})\) be a \hyperref[def:ringed-space]{ringed space}.
  Let \(M\) be a \(\Gamma(X, \sheaf{O}_{X})\) module.
  Consider the \hyperref[def:ringed-space]{morphism of ringed spaces}
  \begin{equation*}
    (\pi, \pi^{\#}): (X, \sheaf{O}_{X}) \to (\{\ast\}, \Gamma(X, \sheaf{O}_{X}))
  \end{equation*}
  
  We define the \textbf{"constant" sheaf} supported at \(M\) as \hyperref[def:inverse-image-ox-mod]{inverse image} of \(M\) via \(\pi\):
  \begin{equation*}
    \sheaf{O}_{X} \otimes_{\Gamma(X, \sheaf{O}_{X})} M \coloneq \pi^{\ast}M = \sheaf{O}_{X} \otimes_{\pi^{-1}\Gamma(X, \sheaf{O}_{X})} \pi^{-1}M = \sheaf{O}_{X} \otimes_{\underline{\Gamma(X, \sheaf{O}_{X})}} \underline{M}
  \end{equation*}

  Note that this defines a \hyperref[def:functor]{functor}
  \begin{equation*}
    \sheaf{O}_{X} \otimes_{\Gamma(X, \sheaf{O}_{X})} - : \Mod(\Gamma(X, \sheaf{O}_{X})) \to \Mod(\sheaf{O}_{X})
  \end{equation*}
\end{definition}

\begin{corollary} \label{cor:global-constant-adjunction-ox-mod}
  Let \((X, \sheaf{O}_{X})\) be a \hyperref[def:ringed-space]{ringed space}.
  The \hyperref[def:global-sections-ox-mod]{global sections functor} and \hyperref[def:constant-sheaf-ox-mod]{constant sheaf functor} of \hyperref[def:sheaf-of-modules]{sheaves of \(\sheaf{O}_{X}\)-modules} are \hyperref[def:adjoint-functors]{adjoint}.

  In formula, if \(\sheaf{F} \in \Mod(\sheaf{O}_{X})\) and \(M \in \Mod(\Gamma(X, \sheaf{O}_{X}))\)
  \begin{equation*}
    \Hom_{\Mod(\sheaf{O}_{X})}(\sheaf{O}_{X} \otimes_{\Gamma(X, \sheaf{O}_{X})} M, \sheaf{F}) \cong \Hom_{\Mod(\Gamma(X, \sheaf{O}_{X}))} (M, \Gamma(X, \sheaf{F}))
  \end{equation*}
\end{corollary}

\begin{definition}[Skyscraper sheaf] \label{def:skyscraper-ox-mod}
  Let \((X, \sheaf{O}_{X})\) be a \hyperref[def:ringed-space]{ringed space} with \(p \in X\).
  Let \(M\) be a \(\sheaf{O}_{X, p}\)-module.
  Consider the \hyperref[def:ringed-space]{morphism of ringed spaces}
  \begin{equation*}
    (\iota_{p}, \iota_{p}^{\#}): (\{\ast\}, \sheaf{O}_{X, p}) \to (X, \sheaf{O}_{X}))
  \end{equation*}
  
  We define the \textbf{skyscraper sheaf} supported at \(M\) as the \hyperref[def:direct-image-ox-mod]{direct image} of \(M\) via \(\iota_{p}\):
  \begin{equation*}
    \iota_{p, \ast}\underline{M} \coloneq \iota_{p, \ast}M
  \end{equation*}

  Note that this defines a \hyperref[def:functor]{functor}
  \begin{equation*}
    \iota_{p, \ast}: \Mod(\sheaf{O}_{X, p}) \to \Mod(\sheaf{O}_{X})
  \end{equation*}
\end{definition}

\begin{definition}[Stalk] \label{def:stalk-ox-mod}
  Let \((X, \sheaf{O}_{X})\) be a \hyperref[def:ringed-space]{ringed space} with \(p \in X\).
  Let \(\sheaf{F}\) be a \hyperref[def:sheaf-of-modules]{sheaf of \(\sheaf{O}_{X}\)-modules}.
  Consider the \hyperref[def:ringed-space]{morphism of ringed spaces}
  \begin{equation*}
    (\iota_{p}, \iota_{p}^{\#}): (\{\ast\}, \sheaf{O}_{X, p}) \to (X, \sheaf{O}_{X}))
  \end{equation*}

  We define the \textbf{stalk} of \(\sheaf{F}\) at \(p\) as \hyperref[def:inverse-image-ox-mod]{inverse image} of \(\sheaf{F}\) via \(\iota_{p}\):
  \begin{equation*}
    \sheaf{F}_{p} \coloneq \iota_{p}^{\ast}\sheaf{F} = \sheaf{O}_{X, p} \otimes_{\iota_{p}^{-1}\sheaf{O}_{X}} \iota_{p}^{-1}\sheaf{F} = \sheaf{O}_{X, p} \otimes_{\sheaf{O}_{X, p}} \sheaf{F}_{p} = \sheaf{F}_{p}
  \end{equation*}
  
  Note that this defines a \hyperref[def:functor]{functor}
  \begin{equation*}
    -_{p} : \Mod(\sheaf{O}_{X}) \to \Mod(\sheaf{O}_{X, p})
  \end{equation*}
\end{definition}

\begin{corollary} \label{cor:stalk-skyscraper-adjunction-ox-mod}
  Let \((X, \sheaf{O}_{X})\) be a \hyperref[def:ringed-space]{ringed space} with \(p \in X\).
  The \hyperref[def:skyscraper-ox-mod]{skyscraper functor} and \hyperref[def:stalk-ox-mod]{stalk functor} of \hyperref[def:sheaf-of-modules]{sheaves of \(\sheaf{O}_{X}\)-modules} are \hyperref[def:adjoint-functors]{adjoint}.

  In formula, if \(\sheaf{F} \in \Mod(\sheaf{O}_{X})\) and \(M \in \Mod(\sheaf{O}_{X, p})\)
  \begin{equation*}
    \Hom_{\Mod(\sheaf{O}_{X, p})}(\sheaf{F}_{p}, M) \cong \Hom_{\Mod(\sheaf{O}_{X})} (\sheaf{F}, \iota_{p, \ast}\underline{M})
  \end{equation*}
\end{corollary}



\begin{lemma}
  Let \((f, f^{\#}): (X, \sheaf{O}_{X}) \to (Y, \sheaf{O}_{Y})\) be a \hyperref[def:ringed-space]{morphism of ringed spaces} and \(p \in X\).
  Let \(\sheaf{G}\) be a \hyperref[def:sheaf-of-modules]{sheaf of \(\sheaf{O}_{Y}\)-modules}.

  As \(\sheaf{O}_{X, p}\) modules we have
  \begin{equation*}
    (f^{\ast}\sheaf{G})_{p} = \sheaf{G}_{f(p)} \otimes_{\sheaf{O}_{Y, f(p)}} \sheaf{O}_{X, p}
  \end{equation*}
\end{lemma}

\end{document}
